\begin{claim}{C.41}{Proved}{A dense operator orbit settles the zero-level kernel}
Let $U$ be a connected admissible surface of C.1, let
$A=\operatorname{Homeo}(U)$ with its compact-open group
topology, and let $H\leq A$ have a dense orbit in $U$.
Then the exact inclusion-induced comparison
\[
 \pi_1(C_A(H),1)\longrightarrow\pi_1(A,1)^H
                                                        \tag{C.41.1}
\]
is injective. In particular, on a hyperbolic surface
$\pi_1(C_{A^0}(H),1)=0$. No local compactness assumption
is imposed on $H$ or on the subgroup generated by
the values of the loop.
\end{claim}

\begin{proof}
\emph{The plane argument.}
Let $f_t$ be a based ambient-null loop of plane
homeomorphisms centralizing a subgroup $H$ with
a dense orbit $D$. Put $G=\langle f_t\rangle$.
Since the loop is based, its maps preserve
orientation, and the orbit results C.28--C.37
apply. Every $h\in H$ transports $G$-orbits
according to the exact identity
\[
                       h(Gx)=G(hx).
                                                        \tag{C.41.2}
\]
If there is no open $G$-orbit, C.34 already
gives the required global equivariant contraction.

Otherwise let $V$ be the nonempty union of the
open $G$-orbits. It is open and $H$-invariant.
Density makes $D$ meet $V$, and invariance then
gives $D\subseteq V$. Consequently $V$ is dense
in the plane.

Fix an open $G$-orbit $O$ and choose $y\in D\cap O$.
Let $H_O=\{h\in H:hO=O\}$ be its operator
stabilizer. The actual orbit equality is
\[
                     H y\cap O=H_O y.
                                                        \tag{C.41.3}
\]
Indeed, $hy\in O$ implies that the two $G$-orbits
$hO=G(hy)$ and $O$ meet, hence are equal. The
reverse inclusion is immediate. Since $Hy=D$
is dense, (C.41.3) is dense in $O$.

Retain the exact restriction homomorphism
$H_O\to\operatorname{Homeo}(O)$, with its
pointwise-fixing kernel. Let $Q_O$ be the
closure of its image in that homeomorphism
group. C.30 proves that $Q_O$ is locally
compact and acts freely and properly on $O$;
C.33 makes it a Lie group of dimension at
most two. Properness makes every $Q_O$-orbit
closed. In particular $Q_Oy$ is a closed
subset of $O$ containing the dense set
$H_Oy$, so
\[
                         Q_Oy=O.
                                                        \tag{C.41.4}
\]
Thus this free proper action is transitive.
Its orbit map $Q_O\to O$ is a homeomorphism,
and $\dim Q_O=2$.

This holds for every open $G$-orbit. In the
notation of C.37, therefore,
\[
 U_2=V,\qquad
 X_2=\overline{U_2}=\R^2.
                                                        \tag{C.41.5}
\]
The contraction proved there is defined on
the whole plane, with jointly continuous
inverses, and centralizes every element of
the original $H$. This settles both alternatives
of the plane argument.

\emph{The actual surface comparison.}
Let a based loop $f_t$ represent a kernel class
in (C.41.1). The inclusion of the fixed
subgroup in $\pi_1(A,1)$ is injective, so
choose an ambient filling
$F:D^2\to A$ of that loop.
For the universal covering
$p:\widetilde U\to U$, every admissible
$\widetilde U$ is a plane. Lift the actual map
\[
 D^2\times\widetilde U\longrightarrow U,\qquad
           (z,x)\longmapsto F(z)(p(x))
                                                        \tag{C.41.6}
\]
with the identity prescribed on the based
boundary fiber. Its domain is simply connected.
As in C.39, lifting the inverse family and
using uniqueness gives an ambient filling by
plane homeomorphisms with jointly continuous
inverses. Its boundary is a closed based
plane loop $\widetilde f_t$.

Let $\widetilde H$ be the group of all lifts
of elements of $H$, including the deck group.
For every $\widetilde h\in\widetilde H$,
the commutator
$[\widetilde f_t,\widetilde h]$ is a deck
transformation depending continuously on $t$
and equal to the identity at $t=0$.
It is therefore always the identity.
The lifted loop centralizes precisely the
same operator lifts required for descent.

Choose $x\in U$ with dense $Hx$ and
$\widetilde x\in p^{-1}(x)$. One has the
exact equality
\[
       \widetilde H\,\widetilde x=p^{-1}(Hx).
                                                        \tag{C.41.7}
\]
One inclusion follows by projection. For the
other, choose a lift of the specified
$h\in H$ and then a deck transformation
taking its value at $\widetilde x$ to the
specified point over $hx$.
The right side is dense because $p$ is open:
every nonempty open set upstairs projects
to a nonempty open set meeting $Hx$.

Apply the plane result to $\widetilde H$.
Its contraction commutes with all deck
transformations and all operator lifts.
The equation
$\overline R_{s,t}p=pR_{s,t}$ therefore
defines a based contraction downstairs
inside $C_A(H)$. The open covering quotient
and its local sections give joint continuity
of these maps and their inverses.
It contracts the original class in the
source of (C.41.1), proving injectivity.
For a hyperbolic $U$, the ambient identity
component is contractible, so its based
centralizer loops all satisfy the kernel
condition. The based contraction stays in
$A^0$.
\end{proof}

\paragraph{Remaining scope.}
C.41 settles every operator action with a
dense orbit. It does not assert that a
general operator group has such an orbit
or supply an equivariant decomposition into
subsets on which this argument applies.
The general operator problem without a dense orbit
in C.12 remains open.
